% ------------------------------------------------------------------
%  Capitolo 10 — Mescolare le carte
%  Parte III
%  Ricerca di riferimento: ricerca 01 §5.3
%  Voci bibliografiche principali: kornell2008, rohrer2007, rohrer2020, brunmair2019
%  Epigrafe: ✔ verificata sul testo (vedi ricerca 03 e registro, ricerca 00)
%  Scheda del capitolo: ricerca/06_indice-ragionato.md, capitolo 10
% ------------------------------------------------------------------
\chapter{Mescolare le carte}
\label{cap:alternanza}

\epigrafe[Dobbiamo imitare le api, come si dice, che vagano e colgono i fiori adatti a fare il miele.]{Apes, ut aiunt, debemus imitari, quae vagantur et flores ad mel faciendum idoneos carpunt.}{Seneca, Epistulae ad Lucilium 84, 3}

\iniziale{Q}{uesto capitolo} \segnaposto{apertura: da scrivere — vedi la scheda nell'indice ragionato}.

\section{Blocchi e mescolanze}

\segnaposto{da scrivere}

\section{Gli esperimenti}

\segnaposto{da scrivere}

\section{Quando funziona e quando no}

\segnaposto{da scrivere}

\section{Come mescolare}

\segnaposto{da scrivere}

\begin{daricordare}
\segnaposto{il principio del capitolo in due righe}
\end{daricordare}

\begin{domande}
  \item \segnaposto{domanda di recupero su questo capitolo}
  \item \segnaposto{domanda di applicazione a una materia che studi}
  \item \segnaposto{domanda di ripresa su un capitolo precedente}
\end{domande}

\begin{sintesi}
  \item \segnaposto{punto di sintesi}
\end{sintesi}

\finecapitolo
